\IfFileExists{stacks-project.cls}{%
\documentclass{stacks-project}
}{%
\documentclass{amsart}
}

% For dealing with references we use the comment environment
\usepackage{verbatim}
\newenvironment{reference}{\comment}{\endcomment}
%\newenvironment{reference}{}{}
\newenvironment{slogan}{\comment}{\endcomment}
\newenvironment{history}{\comment}{\endcomment}

% For commutative diagrams we use Xy-pic
\usepackage[all]{xy}

% We use 2cell for 2-commutative diagrams.
\xyoption{2cell}
\UseAllTwocells

% We use multicol for the list of chapters between chapters
\usepackage{multicol}

% This is generally recommended for better output
\usepackage{lmodern}
\usepackage[T1]{fontenc}

% For cross-file-references
\usepackage{xr-hyper}

% Package for hypertext links:
\usepackage{hyperref}

% For any local file, say "hello.tex" you want to link to please
% use \externaldocument[hello-]{hello}
\externaldocument[introduction-]{introduction}
\externaldocument[conventions-]{conventions}
\externaldocument[sets-]{sets}
\externaldocument[categories-]{categories}
\externaldocument[topology-]{topology}
\externaldocument[sheaves-]{sheaves}
\externaldocument[sites-]{sites}
\externaldocument[stacks-]{stacks}
\externaldocument[fields-]{fields}
\externaldocument[algebra-]{algebra}
\externaldocument[brauer-]{brauer}
\externaldocument[homology-]{homology}
\externaldocument[derived-]{derived}
\externaldocument[simplicial-]{simplicial}
\externaldocument[more-algebra-]{more-algebra}
\externaldocument[smoothing-]{smoothing}
\externaldocument[modules-]{modules}
\externaldocument[sites-modules-]{sites-modules}
\externaldocument[injectives-]{injectives}
\externaldocument[cohomology-]{cohomology}
\externaldocument[sites-cohomology-]{sites-cohomology}
\externaldocument[dga-]{dga}
\externaldocument[dpa-]{dpa}
\externaldocument[sdga-]{sdga}
\externaldocument[hypercovering-]{hypercovering}
\externaldocument[schemes-]{schemes}
\externaldocument[constructions-]{constructions}
\externaldocument[properties-]{properties}
\externaldocument[morphisms-]{morphisms}
\externaldocument[coherent-]{coherent}
\externaldocument[divisors-]{divisors}
\externaldocument[limits-]{limits}
\externaldocument[varieties-]{varieties}
\externaldocument[topologies-]{topologies}
\externaldocument[descent-]{descent}
\externaldocument[perfect-]{perfect}
\externaldocument[more-morphisms-]{more-morphisms}
\externaldocument[flat-]{flat}
\externaldocument[groupoids-]{groupoids}
\externaldocument[more-groupoids-]{more-groupoids}
\externaldocument[etale-]{etale}
\externaldocument[chow-]{chow}
\externaldocument[intersection-]{intersection}
\externaldocument[pic-]{pic}
\externaldocument[weil-]{weil}
\externaldocument[adequate-]{adequate}
\externaldocument[dualizing-]{dualizing}
\externaldocument[duality-]{duality}
\externaldocument[discriminant-]{discriminant}
\externaldocument[derham-]{derham}
\externaldocument[local-cohomology-]{local-cohomology}
\externaldocument[algebraization-]{algebraization}
\externaldocument[curves-]{curves}
\externaldocument[resolve-]{resolve}
\externaldocument[models-]{models}
\externaldocument[functors-]{functors}
\externaldocument[equiv-]{equiv}
\externaldocument[pione-]{pione}
\externaldocument[etale-cohomology-]{etale-cohomology}
\externaldocument[proetale-]{proetale}
\externaldocument[relative-cycles-]{relative-cycles}
\externaldocument[more-etale-]{more-etale}
\externaldocument[trace-]{trace}
\externaldocument[crystalline-]{crystalline}
\externaldocument[spaces-]{spaces}
\externaldocument[spaces-properties-]{spaces-properties}
\externaldocument[spaces-morphisms-]{spaces-morphisms}
\externaldocument[decent-spaces-]{decent-spaces}
\externaldocument[spaces-cohomology-]{spaces-cohomology}
\externaldocument[spaces-limits-]{spaces-limits}
\externaldocument[spaces-divisors-]{spaces-divisors}
\externaldocument[spaces-over-fields-]{spaces-over-fields}
\externaldocument[spaces-topologies-]{spaces-topologies}
\externaldocument[spaces-descent-]{spaces-descent}
\externaldocument[spaces-perfect-]{spaces-perfect}
\externaldocument[spaces-more-morphisms-]{spaces-more-morphisms}
\externaldocument[spaces-flat-]{spaces-flat}
\externaldocument[spaces-groupoids-]{spaces-groupoids}
\externaldocument[spaces-more-groupoids-]{spaces-more-groupoids}
\externaldocument[bootstrap-]{bootstrap}
\externaldocument[spaces-pushouts-]{spaces-pushouts}
\externaldocument[spaces-chow-]{spaces-chow}
\externaldocument[groupoids-quotients-]{groupoids-quotients}
\externaldocument[spaces-more-cohomology-]{spaces-more-cohomology}
\externaldocument[spaces-simplicial-]{spaces-simplicial}
\externaldocument[spaces-duality-]{spaces-duality}
\externaldocument[formal-spaces-]{formal-spaces}
\externaldocument[restricted-]{restricted}
\externaldocument[spaces-resolve-]{spaces-resolve}
\externaldocument[formal-defos-]{formal-defos}
\externaldocument[defos-]{defos}
\externaldocument[cotangent-]{cotangent}
\externaldocument[examples-defos-]{examples-defos}
\externaldocument[algebraic-]{algebraic}
\externaldocument[examples-stacks-]{examples-stacks}
\externaldocument[stacks-sheaves-]{stacks-sheaves}
\externaldocument[criteria-]{criteria}
\externaldocument[artin-]{artin}
\externaldocument[quot-]{quot}
\externaldocument[stacks-properties-]{stacks-properties}
\externaldocument[stacks-morphisms-]{stacks-morphisms}
\externaldocument[stacks-limits-]{stacks-limits}
\externaldocument[stacks-cohomology-]{stacks-cohomology}
\externaldocument[stacks-perfect-]{stacks-perfect}
\externaldocument[stacks-introduction-]{stacks-introduction}
\externaldocument[stacks-more-morphisms-]{stacks-more-morphisms}
\externaldocument[stacks-geometry-]{stacks-geometry}
\externaldocument[moduli-]{moduli}
\externaldocument[moduli-curves-]{moduli-curves}
\externaldocument[examples-]{examples}
\externaldocument[exercises-]{exercises}
\externaldocument[guide-]{guide}
\externaldocument[desirables-]{desirables}
\externaldocument[coding-]{coding}
\externaldocument[obsolete-]{obsolete}
\externaldocument[fdl-]{fdl}
\externaldocument[index-]{index}

% Theorem environments.
%
\theoremstyle{plain}
\newtheorem{theorem}[subsection]{Theorem}
\newtheorem{proposition}[subsection]{Proposition}
\newtheorem{lemma}[subsection]{Lemma}

\theoremstyle{definition}
\newtheorem{definition}[subsection]{Definition}
\newtheorem{example}[subsection]{Example}
\newtheorem{exercise}[subsection]{Exercise}
\newtheorem{situation}[subsection]{Situation}

\theoremstyle{remark}
\newtheorem{remark}[subsection]{Remark}
\newtheorem{remarks}[subsection]{Remarks}

\numberwithin{equation}{subsection}

% Macros
%
\def\lim{\mathop{\mathrm{lim}}\nolimits}
\def\colim{\mathop{\mathrm{colim}}\nolimits}
\def\Spec{\mathop{\mathrm{Spec}}}
\def\Hom{\mathop{\mathrm{Hom}}\nolimits}
\def\Ext{\mathop{\mathrm{Ext}}\nolimits}
\def\SheafHom{\mathop{\mathcal{H}\!\mathit{om}}\nolimits}
\def\SheafExt{\mathop{\mathcal{E}\!\mathit{xt}}\nolimits}
\def\Sch{\mathit{Sch}}
\def\Mor{\mathop{\mathrm{Mor}}\nolimits}
\def\Ob{\mathop{\mathrm{Ob}}\nolimits}
\def\Sh{\mathop{\mathit{Sh}}\nolimits}
\def\NL{\mathop{N\!L}\nolimits}
\def\CH{\mathop{\mathrm{CH}}\nolimits}
\def\proetale{{pro\text{-}\acute{e}tale}}
\def\etale{{\acute{e}tale}}
\def\QCoh{\mathit{QCoh}}
\def\Ker{\mathop{\mathrm{Ker}}}
\def\Im{\mathop{\mathrm{Im}}}
\def\Coker{\mathop{\mathrm{Coker}}}
\def\Coim{\mathop{\mathrm{Coim}}}

% Boxtimes
%
\DeclareMathSymbol{\boxtimes}{\mathbin}{AMSa}{"02}

%
% Macros for moduli stacks/spaces
%
\def\QCohstack{\mathcal{QC}\!\mathit{oh}}
\def\Cohstack{\mathcal{C}\!\mathit{oh}}
\def\Spacesstack{\mathcal{S}\!\mathit{paces}}
\def\Quotfunctor{\mathrm{Quot}}
\def\Hilbfunctor{\mathrm{Hilb}}
\def\Curvesstack{\mathcal{C}\!\mathit{urves}}
\def\Polarizedstack{\mathcal{P}\!\mathit{olarized}}
\def\Complexesstack{\mathcal{C}\!\mathit{omplexes}}
% \Pic is the operator that assigns to X its picard group, usage \Pic(X)
% \Picardstack_{X/B} denotes the Picard stack of X over B
% \Picardfunctor_{X/B} denotes the Picard functor of X over B
\def\Pic{\mathop{\mathrm{Pic}}\nolimits}
\def\Picardstack{\mathcal{P}\!\mathit{ic}}
\def\Picardfunctor{\mathrm{Pic}}
\def\Deformationcategory{\mathcal{D}\!\mathit{ef}}
\usepackage{cleveref}
\crefname{originallemma}{original lemma}{original lemmas}
\Crefname{originallemma}{Original lemma}{Original lemmas}
\makeatletter
\newlabel{etale-cohomology-lemma-category-is-colimit}{{09YU}{source}{}{sourcepointer-1}{}}
\newlabel{etale-cohomology-lemma-category-is-colimit@cref}{{[originallemma][0][]09YU}{[1][0][]source}{}{}{}}
\makeatother
\crefname{originallemma}{original lemma}{original lemmas}
\Crefname{originallemma}{Original lemma}{Original lemmas}
\makeatletter
\newlabel{lemma-base-change-shriek}{{0F7L}{source}{}{sourcepointer-2}{}}
\newlabel{lemma-base-change-shriek@cref}{{[originallemma][0][]0F7L}{[1][0][]source}{}{}{}}
\makeatother
\crefname{originallemma}{original lemma}{original lemmas}
\Crefname{originallemma}{Original lemma}{Original lemmas}
\makeatletter
\newlabel{lemma-finiteness-curves}{{0GKR}{source}{}{sourcepointer-3}{}}
\newlabel{lemma-finiteness-curves@cref}{{[originallemma][0][]0GKR}{[1][0][]source}{}{}{}}
\makeatother
\crefname{originallemma}{original lemma}{original lemmas}
\Crefname{originallemma}{Original lemma}{Original lemmas}
\makeatletter
\newlabel{lemma-loc-constant-shriek-rel-dim-1}{{0GKY}{source}{}{sourcepointer-4}{}}
\newlabel{lemma-loc-constant-shriek-rel-dim-1@cref}{{[originallemma][0][]0GKY}{[1][0][]source}{}{}{}}
\makeatother
\crefname{originallemma}{original lemma}{original lemmas}
\Crefname{originallemma}{Original lemma}{Original lemmas}
\makeatletter
\newlabel{lemma-qf-f-shriek-constructible}{{0GK3}{source}{}{sourcepointer-5}{}}
\newlabel{lemma-qf-f-shriek-constructible@cref}{{[originallemma][0][]0GK3}{[1][0][]source}{}{}{}}
\makeatother
\crefname{originallemma}{original lemma}{original lemmas}
\Crefname{originallemma}{Original lemma}{Original lemmas}
\makeatletter
\newlabel{lemma-relative-triangle-associated-to-open}{{0GKP}{source}{}{sourcepointer-6}{}}
\newlabel{lemma-relative-triangle-associated-to-open@cref}{{[originallemma][0][]0GKP}{[1][0][]source}{}{}{}}
\makeatother
\crefname{originallemma}{original lemma}{original lemmas}
\Crefname{originallemma}{Original lemma}{Original lemmas}
\makeatletter
\newlabel{lemma-stalk-R-f-shriek}{{0GK0}{source}{}{sourcepointer-7}{}}
\newlabel{lemma-stalk-R-f-shriek@cref}{{[originallemma][0][]0GK0}{[1][0][]source}{}{}{}}
\makeatother
\crefname{originallemma}{original lemma}{original lemmas}
\Crefname{originallemma}{Original lemma}{Original lemmas}
\makeatletter
\newlabel{morphisms-lemma-generically-finite}{{02NW}{source}{}{sourcepointer-8}{}}
\newlabel{morphisms-lemma-generically-finite@cref}{{[originallemma][0][]02NW}{[1][0][]source}{}{}{}}
\makeatother
\title{Compact pushforward along the affine line of a torsion constructible sheaf has constructible cohomology}
\author{The Stacks Project Authors}
\date{stacks-project:more-etale.tex:local-snapshot-20261012; GNU FDL 1.2 or later}
\begin{document}
\maketitle
\setcounter{subsection}{1}
\noindent\textit{Editorial note.} All delivered mathematical text is the author's own, copied from the pinned Stacks Project chapter file more-etale.tex; no statement or proof step was rewritten, reordered or removed. The counted span is the author environment and its own proof environment. References inside the span are handed to the assembler as navigation stubs that carry the canonical tag of the cited tag only; whether each of them is load-bearing is recorded on the stub itself.\par\medskip
\noindent\textit{Allowed standard prerequisites (editorial).}\begin{itemize}
\item Constructible sheaves as a colimit (Etale Cohomology, Lemma etale-cohomology-lemma-category-is-colimit, canonical tag 09YU). cited at proof lines 4882-4886. Lets the constructible sheaf F descend along the limit presentation of Y, giving the Noetherian finite dimensional reduction.
\item Relative triangle associated to an open (Lemma lemma-relative-triangle-associated-to-open, canonical tag 0GKP). cited at proof lines 4919-4920. The distinguished triangle for U in the affine line that discards the complementary part Z.
\item Locally constant case (Lemma lemma-loc-constant-shriek-rel-dim-1, canonical tag 0GKY). cited at proof line 4924. The terminal named case: R(U to Y)\_! of the finite type locally constant restriction has constructible cohomology.
\end{itemize}\medskip
\ifdefined\hypertarget\hypertarget{sourcepointer-1}{}\fi
\noindent\textit{Source reference.} Original lemma 09YU: sources/stacks-chow/tags.txt, source lines 12577--12577. Cited at proof lines 4884-4886: F is the pullback of a constructible sheaf on A\textasciicircum{}1 over some finite type Z-subalgebra, giving the Noetherian reduction.\par\medskip
\ifdefined\hypertarget\hypertarget{sourcepointer-2}{}\fi
\noindent\textit{Source reference.} Original lemma 0F7L: sources/stacks-chow/tags.txt, source lines 19011--19011. Cited at proof lines 4906-4907: base change compatibility of Rf\_!F handles the complement in the induction step.\par\medskip
\ifdefined\hypertarget\hypertarget{sourcepointer-3}{}\fi
\noindent\textit{Source reference.} Original lemma 0GKR: sources/stacks-chow/tags.txt, source lines 20696--20696. Cited at proof line 4899: finiteness of H\textasciicircum{}q\_c of the fibre in the base case d = 0.\par\medskip
\ifdefined\hypertarget\hypertarget{sourcepointer-4}{}\fi
\noindent\textit{Source reference.} Original lemma 0GKY: sources/stacks-chow/tags.txt, source lines 20703--20703. Cited at proof line 4924: the terminal case, constructibility of R(U to Y)\_!F for the locally constant restriction.\par\medskip
\ifdefined\hypertarget\hypertarget{sourcepointer-5}{}\fi
\noindent\textit{Source reference.} Original lemma 0GK3: sources/stacks-chow/tags.txt, source lines 20673--20673. Cited at proof lines 4921-4922: the finite complement Z to Y has constructible Rf\_!, discarding it from the triangle.\par\medskip
\ifdefined\hypertarget\hypertarget{sourcepointer-6}{}\fi
\noindent\textit{Source reference.} Original lemma 0GKP: sources/stacks-chow/tags.txt, source lines 20694--20694. Cited at proof lines 4919-4920: the distinguished triangle for the open U in the affine line.\par\medskip
\ifdefined\hypertarget\hypertarget{sourcepointer-7}{}\fi
\noindent\textit{Source reference.} Original lemma 0GK0: sources/stacks-chow/tags.txt, source lines 20670--20670. Cited at proof lines 4896-4898: identifies the stalks of R\textasciicircum{}qf\_!F with the compact cohomology of the geometric fibre in the base case.\par\medskip
\ifdefined\hypertarget\hypertarget{sourcepointer-8}{}\fi
\noindent\textit{Source reference.} Original lemma 02NW: sources/stacks-chow/tags.txt, source lines 3495--3495. Cited at proof lines 4916-4918: after shrinking Y to a dense open, Z to Y is finite.\par\medskip
\begin{lemma}
\label{lemma-constructible-shriek-rel-dim-1}
Let $Y$ be an affine scheme. Let $\Lambda$ be a Noetherian ring.
Let $\mathcal{F}$ be a constructible sheaf of $\Lambda$-modules on
$\mathbf{A}^1_Y$ which is torsion. Then $Rf_!\mathcal{F}$ has constructible
cohomology sheaves where $f : \mathbf{A}^1_Y \to Y$ is the structure morphism.
\end{lemma}

\begin{proof}
Say $\mathcal{F}$ is annihilated by $n > 0$. Then we can replace
$\Lambda$ by $\Lambda/n\Lambda$ without changing $Rf_!\mathcal{F}$.
Thus we may and do assume $\Lambda$ is a torsion ring.

\medskip\noindent
Say $Y = \Spec(R)$. Then, if we write
$R = \bigcup R_i$ as the union of its finite type $\mathbf{Z}$-subalgebras,
we can find an $i$ such that $\mathcal{F}$ is the pullback of
a constructible sheaf of $\Lambda$-modules on $\mathbf{A}^1_{R_i}$, see
\'Etale Cohomology, Lemma \ref{etale-cohomology-lemma-category-is-colimit}.
Hence we may assume $Y$ is a Noetherian scheme of finite dimension.

\medskip\noindent
Assume $Y$ is a Noetherian scheme of finite dimension $d = \dim(Y)$
and $\Lambda$ is torsion. We will prove the result by induction on $d$.

\medskip\noindent
Base case. If $d = 0$, then the only thing to show is that the stalks
of $R^qf_!\mathcal{F}$ are finite $\Lambda$-modules.
If $\overline{y}$ is a geometric point of $Y$, then we have
$(R^qf_!\mathcal{F})_{\overline{y}} = H^q_c(X_{\overline{y}}, \mathcal{F})$
by Lemma \ref{lemma-stalk-R-f-shriek}.
This is a finite $\Lambda$-module by Lemma \ref{lemma-finiteness-curves}.

\medskip\noindent
Induction step. It suffices to find a dense open $V \subset Y$ such
that $Rf_!\mathcal{F}|_V$ has constructible cohomology sheaves. Namely, the
restriction of $Rf_!\mathcal{F}$ to the complement $Y \setminus V$
will have constructible cohomology sheaves by induction
and the fact that formation of $Rf_!\mathcal{F}$ commutes
with all base change (Lemma \ref{lemma-base-change-shriek}).
By definition of constructible sheaves of $\Lambda$-modules, there is a dense
open subscheme $U \subset \mathbf{A}^1_Y$ such that $\mathcal{F}|_U$
is a finite type, locally constant sheaf of $\Lambda$-modules.
Denote $Z \subset \mathbf{A}^1_Y$ the complement (viewed as a reduced
closed subscheme). Note that $U$ contains all the generic points of the
fibres of $\mathbf{A}^1_Y \to Y$ over the generic points
$\xi_1, \ldots, \xi_n$ of the irreducible components of $Y$.
Hence $Z \to Y$ has finite fibres over $\xi_1, \ldots, \xi_n$.
After replacing $Y$ by a dense open (which is allowed), we may
assume $Z \to Y$ is finite, see
Morphisms, Lemma \ref{morphisms-lemma-generically-finite}.
By the distinguished triangle of
Lemma \ref{lemma-relative-triangle-associated-to-open}
and the result for $Z \to Y$ (Lemma \ref{lemma-qf-f-shriek-constructible})
we reduce to showing that
$R(U \to Y)_!\mathcal{F}$ has constructible cohomology sheaves.
This is Lemma \ref{lemma-loc-constant-shriek-rel-dim-1}.
\end{proof}

\end{document}
